% !TEX TS-program = xelatex
% !TEX encoding = UTF-8 Unicode
% !Mode:: "TeX:UTF-8"

\documentclass{resume}
\usepackage{zh_CN-Adobefonts_external}
\usepackage{linespacing_fix}
\usepackage{cite}

\begin{document}
\pagenumbering{gobble}

\name{陈丰逸 \quad Fred}

\basicInfo{
  \email{rei\_utsuho@outlook.com} \textperiodcentered\
  \phone{139~1385~4219} \textperiodcentered\
  \github[UTSUHO]{https://github.com/UTSUHO}
}

%====================
\section{\faCogs\ 技术方向}
前端与平台技术负责人（Front-end \& Platform Tech Lead），长期从事中大型 Web 系统架构设计、工程化建设与交付推进，具备 AI Infra算力平台、科研项目工程化落地经验。
%====================
\section{\faGraduationCap\ 教育背景}

\datedsubsection{\textbf{Rochester Institute of Technology}}{2016.08 -- 2022.01}
\textit{网络与移动计算 B.Sc. in Web \& Mobile Computing}

%====================
\section{\faCogs\ 技术能力}

\begin{itemize}[parsep=0.5ex]
  \item 前端技术：JavaScript  / Vue2 \& Vue3 / Vite / Webpack / SASS / TypeScript / React
  \item 平台与工程化： Git / Docker / Monorepo / CI-CD / 前端架构设计
  \item 系统与可视化：SVG / WebGL / Three.js / GSAP / WebGPU
  \item 后端与数据：Node.js / MySQL / 数据建模 / ETL（Kettle）
  \item 其他：GTK3 / macOS 交叉编译 / AI agent / 需求分析 / 项目规划
\end{itemize}

%====================
% \section{\faUsers\ 工作经历}
% \datedsubsection{\textbf{南京师范大学}}{2024.12 -- 至今}
% \role{研究生就业与职业发展指导顾问（特聘教授聘任）}{兼职}

% \begin{itemize}
%   \item 实际承担研究生就业与职业发展指导工作，面向计算机相关专业及有跨专业就业需求的研究生，提供技术路径规划、工程能力提升与求职能力辅导。
% \end{itemize}
\datedsubsection{\textbf{中科南京信息高铁研究院}}{2022.09 -- 至今}
\role{前端工程师 $\rightarrow$ 前端负责人 $\rightarrow$ 项目技术负责人}{全职}

\begin{itemize}
  \item 承担东数西算工程下南京算力网、AI Infra、社会计算等多项中大型平台项目前端架构设计、技术选型与核心模块实现。
  \item 负责需求拆解、任务分配与工程化指导，在多项目并行与交付周期高度压缩条件下，统筹推进研发进度。
  \item 指导无工程经验研究生开展规范化开发，形成稳定的工程指导与交付机制。
  \item 主导构建体系重构（Webpack $\rightarrow$ Vite），将开发服务器启动时间由约 3 分钟优化至约 20 秒。
  \item 优化关键业务模块性能，将页面加载时间由 21s 降低至 5s。
  \item 在前序项目中完成 WebGL / Three.js / WebGPU 的技术调研与原型验证，并参与多类数据大屏与可视化系统建设，综合经验用于当前项目 3D 渲染与可视化架构设计。
\end{itemize}

\vspace{0.5em}
\texttt{Vue2/3} \textperiodcentered\
\texttt{Node.js} \textperiodcentered\
\texttt{Vite} \textperiodcentered\
\texttt{Webpack} \textperiodcentered\
\texttt{GSAP} \textperiodcentered\
\texttt{Three.js} \textperiodcentered\
\texttt{WEBGPU} \textperiodcentered\


%====================
\section{\faUsers\ 合作与科研项目}

\datedsubsection{\textbf{流浪地球 3 剧组}}{2025 -- 至今}
\role{外部工程技术支持 / 技术顾问}{}

\begin{itemize}
  \item 参与面向长篇叙事文本的多智能体知识图谱构建系统设计与实验方案实施。
  \item 负责论文相关数据集调研与实验复现。
  \item 协助组内开发搭建研发环境与基础工程设施，明确后续技术路线与阶段性目标。
  \item 调研、评估功能需求，引入社区成熟开源项目以支撑系统工程化推进。
\end{itemize}

%====================
\section{\faBook\ 论文成果}

\datedline{\textbf{KGAG}: 基于多代理框架的复杂叙事知识图构建分析器}{2025，ICLR 2026 在投}

\datedline{\textbf{STAGE}: 面向电影剧本的知识图谱、问答与角色扮演基准数据集}{2025，ACL 2026 在投}

%====================
%====================
\section{\faUsers\ 重点项目经历-南研院}
%--------------------

\datedsubsection{\textbf{工热所 CAD 特征识别项目}}{}
\role{项目技术负责人}{}

\begin{itemize}
  \item 以技术负责人角色主导 CAD 相关目标识别 AI 模型项目，主导技术路线制定、核心论文解读与方法论梳理。
  \item 指导研究生完成实验方案设计与实现，持续跟踪并评估实验进展与阶段性结果。
  \item 基于甲方需求、参考论文与阶段目标，明确技术选型与交付边界，围绕 STEP 文件解析与 CAD 目标识别能力开展关键技术调研与可行性论证。
  \item 与项目经理协同评估核心技术实现路径，支撑合同谈判中的技术方案与交付口径制定。
  \item 项目执行阶段统筹研究生推进技术方案落地，协调实验资源与计算环境配置，保障实验按计划推进并满足阶段性交付要求。
\end{itemize}

\vspace{0.3em}
\texttt{CAD} \textperiodcentered\
\texttt{Three.js} \textperiodcentered\
\texttt{OCCT} \textperiodcentered\
\texttt{复现论文实验}
%--------------------

\datedsubsection{\textbf{某科研院所 · 超算云运维管理平台}}{}
\role{前端技术负责人}{}

\begin{itemize}
  \item 作为前端技术负责人，主导基于开源脚手架的中型超算云平台项目建设，负责前端整体技术方案与架构设计。
  \item 承担需求拆解、前端技术选型、复杂模块实现及前端任务分配与进度协调工作。
  \item 在交付周期高度压缩的条件下，同步指导 5 名无工程经验研究生（含部分驻场）开展工程化开发，保障项目按期交付。
  \item 项目交付后持续提供工程开发指导与技术支持，逐步沉淀技术指导机制，支撑平台后续演进。
\end{itemize}

\vspace{0.3em}
\texttt{Vue3} \textperiodcentered\
\texttt{Vite} \textperiodcentered\
\texttt{项目规划}

%--------------------

\datedsubsection{\textbf{某国家级应急信息报送系统}}{}
\role{项目稳定化工程师}{}

\begin{itemize}
  \item 在原项目开发人员在岗期间以关键技术支援身份介入，承担紧急问题处理与核心功能补位，保障项目连续性。
  \item 负责展示与渲染等核心模块实现，确保系统关键功能稳定可用。
  \item 在原开发人员离职后整体接手项目，系统性修复历史缺陷与稳定性问题，显著提升代码可维护性与运行可靠性。
  \item 在维护与迭代过程中指导研究生完成问题定位与调试，推动其形成规范的工程开发与排错流程。
\end{itemize}

\vspace{0.3em}
\texttt{Vue.js} \textperiodcentered\
\texttt{SVG} \textperiodcentered\
\texttt{Flask}

%--------------------

\datedsubsection{\textbf{AI 训推平台（某国家级研究院所 / 上海临港实验室）}}{}
\role{前端开发负责人}{}

\begin{itemize}
  \item 以前端负责人身份主导基于开源“天枢”框架的 AI Infra / 训推平台建设，承担前端整体交付责任。
  \item 负责前端任务拆解与分配，主导高复杂度模块的架构设计与实现，保障核心功能按期落地。
  \item 重构数据集标注模块，将原独立页面改造为基于 dataset type 的抽象化实现，提升代码内聚性并降低维护成本。
  \item 主导使用 Vite 重构原基于 Babel 转译 JSX + Vue CLI 的 Webpack 构建体系，将开发服务器启动时间由约 3 分钟优化至约 20 秒。
  \item 在开源框架图标体系不可扩展的前提下，重新设计并实现项目自用 icon font 打包与引入方案，统一图标管理规范。
  \item 在部门人员变动、工期高度压缩的情况下，合理统筹有限资源，保障项目整体进度与交付质量。
  \item 面向上海临港实验室业务系统，作为系统集成组件参与 AI 训推平台的功能裁剪与定制化集成工作，将既有平台能力适配嵌入甲方现有系统，保障核心功能在新业务场景下的可用性与稳定性。
  \item 直接对接甲方开展需求沟通与使用场景分析，基于实际运行环境提出可行的前端集成方案，独立完成前端适配开发、统一认证系统对接与联调，并配合完成部署与上线，推动问题闭环解决。

\end{itemize}

\vspace{0.3em}
\texttt{JavaScript} \textperiodcentered\
\texttt{Vue.js} \textperiodcentered\
\texttt{Node.js} \textperiodcentered\
\texttt{Vite}


%--------------------

\datedsubsection{\textbf{小易云桌面（云桌面客户端）}}{}
\role{客户端开发 / 工程实现}{}

\begin{itemize}
  \item 参与小易云桌面  客户端的工程实现，面向云桌面 / 虚拟化场景提供本地桌面接入能力。
  \item 基于 GTK3 进行桌面客户端界面开发，完成核心交互界面与基础窗口管理功能实现。
  \item 负责客户端功能模块的工程化落地，包括基础 UI 组件封装、状态展示与用户交互逻辑实现。
  \item 在 OSX/Linux 桌面环境下进行交叉编译windows产物，进行应用调试与问题定位，处理不同发行版环境中的兼容性问题。
  \item 配合整体系统方案，对客户端功能边界与实现复杂度进行评估，支撑产品形态验证与交付。
\end{itemize}

\vspace{0.3em}
\texttt{GTK3} \textperiodcentered\
\texttt{交叉编译} \textperiodcentered\
\texttt{桌面客户端} \textperiodcentered\
\texttt{CSS}

%====================
\section{\faCode\ 个人/实习项目}
%--------------------

\datedsubsection{\textbf{ Isbahn 开发者平台}}{南研院}
\role{前端RD}{}

\begin{itemize}
  \item 与后端同事搭档，利用工作闲暇时间开发工具平台，提高开发效率，为项目推进赋能。
  \item 独立负责前端架构设计、开发与用户体验设计。
  \item 支持动态配置数据源以支撑多项目环境内的应用配置。
  \item 解耦前端配置动态路由表需要后端参与，集成后端代码生成器、SWAGGER文档库等功能。
\end{itemize}

\vspace{0.3em}
\texttt{VUE} \textperiodcentered\
\texttt{ANTDV} \textperiodcentered\
\texttt{vuetify}
%--------------------

\datedsubsection{\textbf{企业级数据仓库系统}}{校内}
\role{数据建模 / ETL 实现}{}

\begin{itemize}
  \item 主导企业级多维数据仓库系统设计与实现，面向多业务部门的复杂数据存储与分析需求。
  \item 负责数据模型设计与事实表、维度表建模工作，梳理业务指标口径并转化为可落地的数据结构。
  \item 使用 Kettle 等 ETL 工具完成数据清洗、转换与加载流程，处理脏数据并保障数据一致性。
  \item 在项目执行过程中承接中途退出成员的核心数据处理工作，独立完成数据清理与导入模块，保障项目按期交付。
\end{itemize}

\vspace{0.3em}
\texttt{数据建模} \textperiodcentered\
\texttt{MySQL} \textperiodcentered\
\texttt{ETL / Kettle} \textperiodcentered\
\texttt{数据仓库}

%--------------------

\datedsubsection{\textbf{法律 ERP 系统（律师事务所）}}{实习}
\role{需求分析 / 项目规划 / 原型设计}{}

\begin{itemize}
  \item 参与律师事务所 ERP 系统的立项与早期规划工作，直接对接业务人员梳理核心业务流程。
  \item 编写 PRD 文档，将业务流程抽象为 ERD 数据模型，并形成可执行的开发文档。
  \item 设计并交付网站与工作流系统原型，用于验证需求合理性与流程可行性。
  \item 协助制定项目代码规范与开发约束，与项目负责人协同解决产品与实施过程中的关键问题。
  \item 开展市场与竞品调研，基于同类产品分析明确系统功能边界与技术发展方向。
\end{itemize}

\vspace{0.3em}
\texttt{需求分析} \textperiodcentered\
\texttt{PRD} \textperiodcentered\
\texttt{数据建模} \textperiodcentered\
\texttt{项目规划}

% \datedsubsection{\textbf{api-disruptor}}{个人}
% \role{Vue.js / express }{}
% \begin{itemize}
%   \item RESTful API 测试工具
%   \item 支持 JSON 格式化输入与输出
%   \item GitHub: https://github.com/UTSUHO/api-disruptor
% \end{itemize}



%====================
\section{\faInfo\ 其他}

\begin{itemize}[parsep=0.5ex]
  \item 语言：中文（母语），英语（熟练听说读写），日语（课程水平 N2）
  \item 兴趣：阅读、旅行、神话研究、古典音乐、电子游戏
\end{itemize}

\end{document}
